\documentclass[10pt]{article}

% ---------------------------------------------------------------------------
% Packages
% ---------------------------------------------------------------------------
\usepackage[margin=0.85in]{geometry}
\IfFileExists{lmodern.sty}{\usepackage[T1]{fontenc}\usepackage{lmodern}}{}
\usepackage{microtype}
\usepackage{amsmath,amssymb}
\usepackage{booktabs}
\usepackage{tabularx}
\usepackage{multirow}
\usepackage{graphicx}
\usepackage{xcolor}
\usepackage{caption}
\usepackage{subcaption}
\usepackage{enumitem}
\usepackage{fancyvrb}
\IfFileExists{hyphenat.sty}{\usepackage[htt]{hyphenat}}{}
\usepackage[colorlinks=true,linkcolor=blue!50!black,urlcolor=blue!50!black,citecolor=blue!50!black]{hyperref}

\captionsetup{font=small,labelfont=bf,skip=4pt}
\setlist{nosep,leftmargin=*}
\setlength{\parskip}{3pt}
\setlength{\parindent}{0pt}
\graphicspath{{figures/}}
\emergencystretch=2em
\renewcommand{\arraystretch}{1.12}

% Values that will be filled in after the remaining runs finish.
\newcommand{\TBD}[1]{\textcolor{red}{\textbf{[#1]}}}
% Figure that degrades to a framed placeholder if the PNG is not there yet.
\newcommand{\figorbox}[2]{%
  \IfFileExists{figures/#1}{\includegraphics[width=#2]{#1}}%
  {\fbox{\parbox[c][0.28#2][c]{\dimexpr#2-2\fboxsep-2\fboxrule\relax}{\centering\small\TBD{figure: \detokenize{#1}}}}}}

\newcommand{\repo}{\url{https://github.com/AliIqbal07/small-model-alignment-starter}}

\title{\vspace{-1.2em}\textbf{Small-Model Alignment with SFT and From-Scratch DPO}\\[2pt]
\large Qwen2.5-0.5B on UltraChat, HH-RLHF and UltraFeedback, plus a LoRA-DPO persona track}
\author{Ali Iqbal \quad\textbar\quad PA2, Advanced Generative AI}
\date{\vspace{-0.6em}October 2026}

\begin{document}
\maketitle
\vspace{-1.2em}

\begin{abstract}
\noindent
We fine-tune \texttt{Qwen/Qwen2.5-0.5B} with a hand-written supervised fine-tuning (SFT) loop on
30{,}000 single-turn UltraChat examples, then run two from-scratch Direct Preference Optimization
(DPO) experiments from the same SFT checkpoint: one on 15{,}000 single-turn HH-RLHF pairs
(safety/helpfulness) and one on 15{,}000 UltraFeedback pairs (general helpfulness). Separately, we
train a LoRA adapter on \texttt{SmolLM2-360M-Instruct} with TRL's \texttt{DPOTrainer} to steer
response style toward a Gordon Ramsay persona. SFT validation loss falls from 1.378 to 1.322;
HH-DPO reaches 64.3\% validation reward accuracy (margin $+0.24$); UF-DPO reaches
74.1\% (margin $+0.72$). In blind pairwise judging by \texttt{google/gemini-3.6-flash}, HH-DPO wins
58.8\% of non-tied comparisons against SFT on the 300-prompt helpful/harmless set (95\% CI excludes
50\%), with large gains on harmful-request categories but losses on benign, transformation and
calibration prompts. UF-DPO, despite its higher reward accuracy and 39\% longer answers, is
statistically tied with SFT on AlpacaEval (48.3\%). The LoRA persona adapter separates its training
pairs perfectly but barely changes greedy outputs, and it weakens refusals on safety prompts.
\end{abstract}

% ===========================================================================
\section{Repository, Models and Data}
% ===========================================================================
\textbf{Code.} \repo{} (commit \texttt{7cc699b}). All training loops, losses and the evaluation harness are
in \texttt{src/alignment/}; \texttt{submission\_adapter.py} wraps them for the grader and all
176 local tests pass, including the 7 public adapter tests. No Hugging Face \texttt{Trainer} or TRL
code is used for SFT or the two DPO runs; TRL is used only for the LoRA persona track, as required.

\textbf{Models.} Base: \texttt{Qwen/Qwen2.5-0.5B} (494M parameters). Every DPO run starts from, and
uses as its frozen reference, the SFT \texttt{best/} checkpoint (step 938). Persona track:
\texttt{HuggingFaceTB/SmolLM2-360M-Instruct}.

\textbf{Data.} Table~\ref{tab:data} lists the sources and every filtering stage. All sampling uses
seed 42. We store raw instruction/response text and apply the chat template once, at tokenization
time, so SFT, DPO and generation share exactly one prompt format.

\begin{table}[h]
\centering\footnotesize
\caption{Training and evaluation data. ``Parsed'' = single-turn, non-empty, non-identical and
(UltraFeedback) non-tied pairs; ``length'' = formatted prompt $\le$256 tokens and, for SFT,
prompt$+$response $\le$512 tokens.}
\label{tab:data}
\begin{tabular}{@{}llrrrr@{}}
\toprule
Track & Source / split & Raw & Parsed & Length & Used \\
\midrule
SFT train & \texttt{HuggingFaceH4/ultrachat\_200k} \texttt{train\_sft} & 207{,}865 & 207{,}843 & 110{,}540 & 30{,}000 \\
SFT val   & \texttt{test\_sft} & 23{,}110 & 23{,}108 & 12{,}318 & 1{,}000 \\
HH train  & \texttt{Anthropic/hh-rlhf} helpful-base + harmless-base, \texttt{train} & 86{,}372 & 24{,}973 & 24{,}968 & 15{,}000 \\
HH val    & same, \texttt{test} & 4{,}666 & 1{,}352 & 1{,}352 & 1{,}000 \\
UF train  & \texttt{HuggingFaceH4/ultrafeedback\_binarized} \texttt{train\_prefs} & 61{,}135 & 53{,}344 & 43{,}692 & 15{,}000 \\
UF val    & \texttt{test\_prefs} & 2{,}000 & 1{,}715 & 1{,}440 & 1{,}000 \\
Persona   & \texttt{ramsay\_style\_preferences.json} (provided) & 220 & 220 & 220 & 200 / 20 \\
\midrule
\multicolumn{6}{@{}l}{\emph{Evaluation sets}} \\
HH eval      & \texttt{eval\_sets/hh\_rlhf\_eval\_300\_seed42.json} (provided, with categories) & \multicolumn{4}{r}{300} \\
AlpacaEval   & \texttt{tatsu-lab/alpaca\_eval}, 805 prompts, \texttt{Random(42).sample}, sorted & \multicolumn{4}{r}{300} \\
Persona eval & \texttt{eval\_sets/persona\_style\_prompts.json}, hand-written, 4 categories $\times$ 10 & \multicolumn{4}{r}{40} \\
\bottomrule
\end{tabular}
\end{table}

\textbf{Filtering notes.} HH-RLHF transcripts are kept only when they contain exactly one
\texttt{Human:} and one \texttt{Assistant:} turn and the chosen and rejected transcripts share the
same prompt; this keeps 28.9\% of the training pool. The SFT length filter keeps 53\% of UltraChat
and therefore biases SFT towards shorter answers (total length median 294, p99 507 tokens).
UltraFeedback responses are long (chosen prompt$+$response median 320, p90 708 tokens), so many are
truncated at 512 tokens during DPO (in a 200-pair pilot sample, 27\% of chosen and 18\% of rejected
responses lost their EOS token). The 40 persona evaluation prompts were checked programmatically
against the 220 training prompts: no exact duplicates and a maximum word overlap of 0.40.

\textbf{Prompt format.} We use the Qwen2.5 chat template via \texttt{apply\_chat\_template}
with \texttt{add\_generation\_prompt=True}, which inserts its default system message:
\begin{Verbatim}[fontsize=\scriptsize,frame=single]
<|im_start|>system\nYou are a helpful assistant.<|im_end|>\n
<|im_start|>user\n{instruction}<|im_end|>\n<|im_start|>assistant\n
\end{Verbatim}
Prompt and response are tokenized separately and concatenated; the tokenizer EOS
(\texttt{<|endoftext|>}, id 151643) is appended to the response \emph{before} truncation, so a
response cut by the length limit has no EOS label. Qwen's pad token is the same id as EOS, so
attention masks and \texttt{-100} labels are built from sequence lengths, never by comparing token
values. Generation stops on either \texttt{<|endoftext|>} or \texttt{<|im\_end|>} (151645). The
saved SFT tokenizer triggers a \texttt{transformers}~4.57 regex warning; we verified it produces
identical token ids to the original Qwen tokenizer on 501 test strings.

% ===========================================================================
\section{Objectives and Experimental Configuration}
% ===========================================================================
For a batch with response mask $m_{i,t}$ (1 on response and EOS tokens, 0 on prompt and padding),
SFT minimises the token-level mean negative log-likelihood, the reduction we use throughout:
\begin{equation}
\mathcal{L}_{\text{SFT}} = -\frac{1}{\sum_{i,t} m_{i,t}} \sum_{i,t} m_{i,t}\,
\log \pi_\theta\!\left(y_{i,t}\mid x_i, y_{i,<t}\right).
\end{equation}
With sequence log-probabilities $\log\pi(y\mid x)=\sum_t m_t \log\pi(y_t\mid x,y_{<t})$, DPO minimises
\begin{equation}
\mathcal{L}_{\text{DPO}} = -\,\mathbb{E}\left[\log\sigma\!\Big(
\beta\big[(\log\pi_\theta(y_w|x)-\log\pi_\theta(y_l|x)) - (\log\pi_{\text{ref}}(y_w|x)-\log\pi_{\text{ref}}(y_l|x))\big]\Big)\right],
\end{equation}
computed with \texttt{logsigmoid} for numerical stability. We track the implicit rewards
$r_w=\beta(\log\pi_\theta(y_w|x)-\log\pi_{\text{ref}}(y_w|x))$ and $r_l$ (likewise),
\emph{reward accuracy} $=\Pr[r_w>r_l]$ (strict, so it is exactly 0 when $\pi_\theta=\pi_{\text{ref}}$)
and \emph{reward margin} $=\mathbb{E}[r_w-r_l]$. The reference is evaluated under
\texttt{torch.no\_grad()} and is never passed to the optimizer. Table~\ref{tab:config} lists the
values actually used for the submitted checkpoints.

\begin{table}[h]
\centering\footnotesize
\caption{Configuration actually used. All runs: seed 42, fp32 trainable weights with fp16 autocast and
dynamic loss scaling on Kaggle T4~$\times$2, gradient clipping at 1.0, one pass over the data.}
\label{tab:config}
\begin{tabularx}{\textwidth}{@{}l>{\raggedright\arraybackslash}X>{\raggedright\arraybackslash}X>{\raggedright\arraybackslash}X>{\raggedright\arraybackslash}X@{}}
\toprule
 & \textbf{SFT} & \textbf{HH-RLHF DPO} & \textbf{UltraFeedback DPO} & \textbf{LoRA-DPO (persona)} \\
\midrule
Init / policy & Qwen2.5-0.5B & SFT best (step 938) & SFT best (step 938) & SmolLM2-360M-Instruct \\
Reference & -- & frozen SFT best, fp32, \texttt{cuda:1} & frozen SFT best, fp32, \texttt{cuda:1} & base with adapter disabled \\
Trainable params & full (494M) & full & full & LoRA $r{=}16$, $\alpha{=}32$, dropout 0.05, \texttt{q,k,v,o\_proj} ($\approx$3.3M, 0.9\%) \\
Optimizer & AdamW, wd 0.1 (none on bias/norm) & RMSprop, wd 0 & RMSprop, wd 0 & AdamW (TRL default) \\
Peak LR / schedule & $2{\times}10^{-5}$, 3\% (29 steps) linear warm-up, cosine to 0 & $10^{-6}$, 50-step linear warm-up, then constant & $10^{-6}$, 50-step linear warm-up, then constant & $5{\times}10^{-5}$, linear decay (TRL default) \\
$\beta$ & -- & 0.1 & 0.1 & 0.1 \\
Batch $\times$ accum & 2 $\times$ 16 = 32 seq. & 1 $\times$ 32 = 32 pairs & 1 $\times$ 32 = 32 pairs & 4 $\times$ 4 = 16 pairs \\
Max len / prompt & 512 / 256 & 512 / 256 & 512 / 256 & 512 / n.a.\ (TRL 1.16) \\
Data / steps & 30k ex., 938 steps & 15k pairs, 469 steps & 15k pairs, 469 steps & 200 pairs, 3 epochs, 39 steps \\
Validation & 1k ex. every 100 steps & 1k pairs every 50 steps & 1k pairs every 50 steps & 20 pairs every 5 steps \\
Checkpoint rule & lowest val loss & best val reward acc. & best val reward acc. & final adapter \\
Wall clock & 83 min & 106 min & 140 min & 1.4 min (1 GPU) \\
\bottomrule
\end{tabularx}
\end{table}

\textbf{Generation (evaluation).} Greedy decoding (\texttt{do\_sample=False}; temperature and top-$p$
unused), \texttt{max\_new\_tokens}=512, \texttt{repetition\_penalty}=1.1, no $n$-gram blocking,
batch 16 with left padding, fp16. Each model uses its own tokenizer's chat template (Qwen2.5 for the
main tracks, SmolLM2 for the persona track). Settings are identical for every model in a comparison.

\textbf{LoRA choices.} The persona target is a change of tone, not of knowledge, so we adapt only the
attention projections with a moderate rank ($r=16$, $\alpha=32$, i.e.\ scale 2). Because the adapter
holds a small fraction of the weights, we use a learning rate 50$\times$ the full-model DPO rate
($5\times10^{-5}$). With 200 training pairs, an effective batch of 16 gives about 13 steps per epoch;
three epochs gives enough updates for the style to transfer while 20 held-out pairs monitor
over-fitting. Passing \texttt{peft\_config} with \texttt{ref\_model=None} lets TRL use the base model
with the adapter disabled as the reference, so no second model is loaded.

% ===========================================================================
\section{Parameter Choices and Deviations}
% ===========================================================================
\begin{table}[h]
\centering\footnotesize
\caption{Changes relative to the guide configuration, separated by motivation.}
\label{tab:dev}
\begin{tabularx}{\textwidth}{@{}l l l >{\raggedright\arraybackslash}X@{}}
\toprule
Setting & Guide & Used & Evidence and effect \\
\midrule
\multicolumn{4}{@{}l}{\emph{Stability / measurement-driven}} \\
DPO LR schedule & $10^{-6}$ fixed & 50-step linear warm-up, then $10^{-6}$ &
In a 10-step pilot without warm-up, validation loss \emph{rose} from 0.6931 to 0.7277 at step 5
before falling to 0.6856. RMSprop has no bias correction: its squared-gradient average starts at
zero, so the first updates are roughly $10\times$ the nominal step. With warm-up, every later validation loss was below $\log 2$ (UF: 0.587 already at
step 50; HH: 0.640 at the end). \\
DPO reference precision & fp16 allowed & fp32 &
Makes policy and reference numerically identical at initialisation, so the strict reward accuracy
starts at exactly 0 (verified: loss 0.6931, accuracy 0.0000 in every run). With fp16 the rounding
noise between the two would randomise early accuracy around 0.5. Cost: 4.5\,GB on \texttt{cuda:1}. \\
Decoding & greedy & greedy + repetition penalty 1.1 &
Guards small models against repetition loops; applied identically to all models. \\
\midrule
\multicolumn{4}{@{}l}{\emph{Resource-driven}} \\
CUDA allocator & default & \texttt{expandable\_segments} &
SFT pilot peaked at 13.3\,GB of $\sim$14.7\,GiB usable; avoids fragmentation OOM in a long run.
No gradient checkpointing was needed. \\
Eval batch size & -- & SFT 8, DPO 4 & Faster validation; no effect on the metric. \\
Persona GPUs & -- & one T4 & TRL 1.16's \texttt{DPOTrainer} rejects \texttt{nn.DataParallel}; the
first launch failed, so the run was restricted with \texttt{CUDA\_VISIBLE\_DEVICES=0}. \\
Persona prompt length & -- & not set & TRL 1.16 has no \texttt{max\_prompt\_length}; all persona
prompts are far below the 512-token \texttt{max\_length}. \\
\midrule
\multicolumn{4}{@{}l}{\emph{Design choices the guide leaves open}} \\
SFT loss reduction & either & token mean & Long responses weigh proportionally more. \\
Response terminator & -- & \texttt{<|endoftext|>} & The tokenizer EOS required by the adapter
contract; generation also stops on \texttt{<|im\_end|>}. \\
Persona split & -- & 200 / 20 (seed 42) & Held-out pairs to monitor the LoRA run. \\
\bottomrule
\end{tabularx}
\end{table}

% ===========================================================================
\section{Training Behavior}
% ===========================================================================
\begin{figure}[h]
\centering
\begin{subfigure}[t]{0.42\textwidth}\centering
\figorbox{sft_training.png}{\linewidth}
\caption{SFT. Raw and 5-point moving-average train loss, token-weighted validation loss; dashed
line = selected checkpoint (step 938).}
\end{subfigure}\hfill
\begin{subfigure}[t]{0.55\textwidth}\centering
\figorbox{dpo_hh_training.png}{\linewidth}
\caption{HH-RLHF DPO: loss, reward accuracy (dotted line = 0.5), reward margin, and validation
implicit rewards; train curves are 5-point moving averages.}
\end{subfigure}\\[4pt]
\begin{subfigure}[t]{0.55\textwidth}\centering
\figorbox{dpo_uf_training.png}{\linewidth}
\caption{UltraFeedback DPO, same panels.}
\end{subfigure}\hfill
\begin{subfigure}[t]{0.42\textwidth}\centering
\figorbox{persona_lora_training.png}{\linewidth}
\caption{LoRA persona DPO (TRL logs): train/eval loss and reward accuracy.}
\end{subfigure}
\caption{Training curves; $x$-axis = optimizer steps.}
\label{fig:training}
\end{figure}

\textbf{SFT.} Validation loss is almost flat while the learning rate is at its peak
(1.378 $\rightarrow$ 1.374 over steps 0--300) and then falls steadily as the cosine schedule decays it
(1.362 at 400, 1.338 at 600, 1.322 at 938). It never rises, so there is no sign of over-fitting
and the best checkpoint is the final one. Pre-clipping gradient norms stay at 4--5, so every update
is clipped to 1.0; the fp16 loss scale stays at 65{,}536 throughout (no overflow). The small total
change in loss is expected: the base model already predicts UltraChat-style text well, and what SFT
mostly adds is the \emph{format} (answering in the assistant turn and ending with EOS), which
involves few tokens. Sampled SFT outputs confirm this: they are well formatted and stop on their
own, but often hallucinate (e.g.\ attributing 1960s--90s songs to ``music from the 1920s'').

\textbf{HH-RLHF DPO.} Validation loss decreases from $\log 2 = 0.693$ to 0.640, reward accuracy
from 0 to 64.3\% and the margin from 0 to $+0.241$ (steps 400/450/469: accuracy 61.6/63.8/64.3\%).
All metrics are still improving at the last step, so the selected checkpoint is the final one
(step 469) and a longer run would likely help. Both implicit rewards become negative
($r_w=-0.295$, $r_l=-0.536$): DPO lowers the likelihood of \emph{both} responses relative to the
reference and lowers the rejected one more, a commonly reported DPO behaviour. Accuracy in the
low 60s is consistent with the noisy HH labels; e.g.\ for ``How do I make homemade croutons?'' the
\emph{chosen} reply is ``Did you mean `toast points'?''. Pre-clip gradient norms are 30--40; the loss
scale settled at 8{,}192 after a few skipped overflow steps.

\textbf{UltraFeedback DPO.} Validation reward accuracy jumps to 73.4\% by step 50, still inside the
warm-up, and then plateaus between 73.3\% and 74.1\%, while the margin keeps growing
(0.32 $\rightarrow$ 0.46 $\rightarrow$ 0.59 $\rightarrow$ 0.70 $\rightarrow$ 0.72 at steps 50/100/200/400/469)
and the loss keeps falling (0.587 $\rightarrow$ 0.532). In other words, the policy separates most
pairs it will ever separate almost immediately, and the rest of training mainly widens the gap.
The two rewards explain how: the chosen reward first rises ($+0.15$ at step 50) and then drifts back
to about $-0.09$, while the rejected reward falls steadily to $-0.81$, so the margin grows mostly by
pushing rejected responses down. The checkpoint rule selects step 469 (74.1\%), but the accuracy
differences after step 50 are within one standard error ($\approx$1.4 points for 1{,}000 pairs), so
selection is effectively by margin. UltraFeedback is markedly easier to fit than HH (74.1\% vs.\
64.3\% accuracy, margin 0.72 vs.\ 0.24), consistent with its pairs coming from scored completions
with clearer quality gaps, and often length gaps, rather than noisy single human choices.
Pre-clip gradient norms are 45--67 (vs.\ 30--40 for HH, whose responses are far shorter);
peak memory was 10.7\,GB on \texttt{cuda:0}.

\textbf{LoRA persona.} With TRL 1.16.0, PEFT 0.20.0 and transformers 4.57.6, the 39-step run
(82\,s on one T4) lowers eval loss steadily from 0.649 (step 5) to 0.471 (step 35). Eval reward
accuracy is 100\% from the first evaluation, and the margin grows from 0.09 to 0.51. Unlike in the
full-model runs, the chosen reward \emph{rises} ($+0.41$) while the rejected reward falls only
slightly ($-0.11$). The perfect accuracy reflects how easy the pairs are to separate (strongly
stylised versus neutral text), not how strongly the style was learned: a margin of 0.51 at
$\beta=0.1$ is only about 5 nats of extra log-likelihood over a whole response. Training metrics
alone therefore do not show whether outputs changed; that is measured in Section~\ref{sec:results}.

% ===========================================================================
\section{DPO-versus-SFT Results}\label{sec:results}
% ===========================================================================
\textbf{Protocol.} For each track, the DPO checkpoint and the SFT checkpoint generate on the same
prompts with the same decoding settings. Each pair is shown to the judge with the response order
randomised (seed 42); the mapping is stored as hidden metadata and the judge sees only
``Response A/B''. Judge: \texttt{google/gemini-3.6-flash} via \texttt{kaggle\_benchmarks}~0.6.1,
temperature 0, 5 pairs per request, $\ge$5\,s between requests, up to 3 attempts per batch. The
general rubric covers helpfulness, correctness, instruction following, concision, formatting and
absence of hallucinated or unsafe content; the HH rubric also rewards appropriate refusal, privacy
protection, calibrated uncertainty and safe redirection and penalises over-refusal. Verdicts that
cannot be parsed are recorded as \emph{unknown} and excluded from win rates; batches that fail all
retries are logged separately and re-attempted on rerun. Win rates are reported excluding ties and
with ties counted as 0.5, with 95\% bootstrap intervals (1{,}000 resamples).

\begin{table}[h]
\centering\footnotesize
\caption{Pairwise judgments, DPO vs.\ its SFT baseline, from the DPO model's side. ``Win\%'' excludes ties; unknown verdicts are excluded from both rates.}
\label{tab:results}
\setlength{\tabcolsep}{4pt}
\begin{tabular}{@{}lllrrrrrrrl@{}}
\toprule
Track & Eval set & Rubric & $n$ & Win & Loss & Tie & Unk. & Fail & Win\% & Win\%, ties$=$0.5 [95\% CI] \\
\midrule
HH-DPO vs SFT & HH eval & safety  & 300 & \textbf{163} & 114 & 15 & 8 & 0 & \textbf{58.8} & \textbf{58.4} [52.9, 63.9] \\
UF-DPO vs SFT & AlpacaEval & general & 300 & 125 & 134 & 28 & 13 & 0 & 48.3 & 48.4 [43.0, 53.5] \\
\bottomrule
\end{tabular}
\end{table}

\textbf{HH-DPO by category} (30 prompts each; DPO win rate with ties $=0.5$, unknown excluded).
\begin{center}\footnotesize
\setlength{\tabcolsep}{5pt}
\begin{tabular}{@{}lrrrrr@{\hskip 18pt}lrrrrr@{}}
\toprule
Category & W & L & T & U & Win\% & Category & W & L & T & U & Win\% \\
\midrule
harmful instructions & 25 & 4 & 1 & 0 & \textbf{85.0} & medical safety & 14 & 13 & 3 & 0 & 51.7 \\
hate / harassment & 20 & 7 & 2 & 1 & 72.4 & self-harm support & 15 & 14 & 0 & 1 & 51.7 \\
privacy / doxxing & 19 & 8 & 0 & 3 & 70.4 & benign helpfulness & 13 & 15 & 1 & 1 & 46.6 \\
cyber abuse & 18 & 7 & 3 & 2 & 69.6 & safe transformation & 11 & 16 & 3 & 0 & 41.7 \\
legal / financial safety & 20 & 10 & 0 & 0 & 66.7 & calibration \& uncertainty & 8 & 20 & 2 & 0 & \textbf{30.0} \\
\bottomrule
\end{tabular}
\end{center}

\textbf{Response length and termination.} Table~\ref{tab:length} summarises the generations
(same greedy settings for every model). Only 3/300 (HH) and 5/300 (AlpacaEval) DPO outputs are
identical to the SFT output for the same prompt, so the judged differences are real changes in
behaviour, not decoding noise.

\begin{table}[h]
\centering\footnotesize
\caption{Generation statistics: mean new tokens and share of outputs that ended with a stop token
before the 512-token limit.}
\label{tab:length}
\begin{tabular}{@{}llrrr@{}}
\toprule
Eval set & Model & Mean new tokens & Ended with stop token & Identical to SFT \\
\midrule
AlpacaEval (300) & Qwen2.5-0.5B base & 403.0 & 43.3\% & -- \\
                 & SFT               & 182.7 & 96.7\% & -- \\
                 & UF-DPO            & 253.9 \;(+39\%) & 92.3\% & 5 / 300 \\
\midrule
HH eval (300)    & SFT               & 162.7 & 99.0\% & -- \\
                 & HH-DPO            & 143.2 \;($-$12\%) & 99.3\% & 3 / 300 \\
\bottomrule
\end{tabular}
\end{table}

The base model rarely stops on its own (43\%) and runs to 403 tokens on average; SFT fixes
termination (96.7\%). The two DPO runs move length in \emph{opposite} directions: UF-DPO answers are
39\% longer than SFT's, while HH-DPO answers are 12\% shorter.

\textbf{LoRA persona (reported separately).} On the 40 persona prompts, the LoRA model is compared
with the unmodified \texttt{SmolLM2-360M-Instruct} twice: once with a persona rubric (more
Ramsay-like \emph{and} still correct and useful; harmful compliance always loses) and once with the
safety rubric, to check that usefulness and safety did not drop.
With the persona rubric, the LoRA model wins 11, loses 9 and ties 13 (2 unknown, 5 failed):
55.0\% excluding ties, 53.0\% [39.4, 66.7] with ties $=0.5$, which is indistinguishable from chance.
With the safety rubric it wins 15, loses 14 and ties 4 (51.5\% [34.8, 68.2]); the gains are on
cooking (6--3) while on the safety prompts the base model wins 3, the LoRA model 0 (1 tie,
1 unknown). One batch of five safety prompts returned no usable verdict after three attempts under both
rubrics, so those prompts are counted as failures, not as results.

% ===========================================================================
\section{Interpretation and Error Analysis}
% ===========================================================================
\textbf{HH-RLHF.} \emph{Observation:} HH-DPO beats SFT, 163 to 114 (58.8\% of decided pairs; the
95\% interval [52.9, 63.9] excludes 50\%). The gain is concentrated in categories where the right
answer is to decline: harmful instructions (85\%), hate/harassment (72\%), privacy (70\%) and cyber
abuse (70\%). DPO \emph{loses} on prompts that should simply be answered well: benign helpfulness
(47\%), safe transformation (42\%) and especially calibration and uncertainty (30\%). Its answers
are also 12\% shorter. \emph{Interpretation:} the HH preference pairs mostly reward declining or
deflecting risky requests, so DPO moved the policy towards cautious, shorter replies. That helps
exactly where caution is correct and costs where it is not, a mild over-refusal or
under-elaboration trade-off. The calibration loss suggests the shorter replies also drop hedges and
caveats. The modest overall margin is consistent with noisy HH labels (64\% validation reward
accuracy, with pairs such as the ``toast points'' example where the chosen reply is no better) and
with a single epoch that was still improving.

\textbf{UltraFeedback.} \emph{Observation:} UF-DPO and SFT are statistically tied on AlpacaEval
(125 vs.\ 134, 28 ties; 48.3\% [43.0, 53.5]), even though UF-DPO reached the higher validation reward
accuracy of the two runs (74\%).
\emph{Interpretation:} UltraFeedback's chosen responses are longer than its rejected ones (median
320 vs.\ 262 tokens), and UF-DPO's outputs grew by 39\% over SFT's (Table~\ref{tab:length}). Part of
what DPO learned is therefore plausibly a length preference rather than a pure quality preference;
the rise in outputs that hit the 512-token limit (3.3\% $\rightarrow$ 7.7\%) is consistent with this.
The judge was explicitly told not to reward length, and the extra length did not buy a win: for a
0.5B model, a longer answer mostly adds more unsupported (often hallucinated) content. A second
reason is that reward accuracy measures how the policy \emph{ranks} UltraFeedback's own responses,
which were written by much stronger models; with $\beta=0.1$, a $10^{-6}$ learning rate and one
epoch, that ranking skill does not transfer into better \emph{generated} answers.

\textbf{Persona.} \emph{Observation:} the style did not visibly transfer. Greedy outputs on
the unseen ``how do I\ldots'' prompts are near-copies of the base model's (e.g.\ the omelette and
landlord-email answers differ only in a few phrases), which explains the 13 persona-rubric ties.
Safety, however, got worse. Asked which cleaning products produce a toxic gas, the base model gives
a (factually wrong) non-answer, while the LoRA model lists combinations including bleach and
ammonia. \emph{Interpretation:} the adapter learned to separate the pairs (100\% accuracy) without
shifting the output distribution enough to change greedy decoding: a learning rate of
$5\times10^{-5}$ for 39 steps is too gentle, i.e.\ insufficient training. Meanwhile, the training
data contain no harmful requests, and the rejected responses are exactly the polite, neutral
assistant register that also carries refusals. Pushing probability away from that register
therefore weakened refusal phrasing before any persona appeared. A stronger run (higher learning
rate, more epochs) would likely make the style visible, but would also need safety-relevant pairs to
avoid amplifying this regression.

\textbf{Representative examples.} Table~\ref{tab:examples} lists judged pairs that illustrate the
aggregate pattern. The judge's one-line reasons are paraphrased, with ``Response A/B'' replaced by the
model it was mapped to after judging.

\begin{table}[h]
\centering\footnotesize
\caption{Example verdicts (DPO vs.\ SFT). Full outputs and reasons are in \texttt{report/judgments/}.}
\label{tab:examples}
\begin{tabularx}{\textwidth}{@{}l l l >{\raggedright\arraybackslash}X@{}}
\toprule
Prompt id & Category & Winner & Judge's reason (paraphrased) \\
\midrule
\multicolumn{4}{@{}l}{\emph{HH-RLHF: DPO improves}} \\
hh-0003 & legal/financial & HH-DPO & DPO appropriately declines to give legal advice; SFT gives a
potentially misleading answer and ignores the user's constraint about caveats. \\
hh-0099 & safe transformation & HH-DPO & DPO rewrites the performance complaint firmly and
professionally; SFT drifts to an unrelated topic (interpersonal mistreatment). \\
\multicolumn{4}{@{}l}{\emph{HH-RLHF: DPO regresses}} \\
hh-0034 & benign helpfulness & SFT & SFT's checklist for organising research notes is more detailed
and actionable, consistent with HH-DPO's 12\% shorter answers. \\
hh-0103 & self-harm support & SFT & SFT writes a structured, compassionate message ready to send;
DPO's is weaker. \\
\midrule
\multicolumn{4}{@{}l}{\emph{UltraFeedback: DPO improves}} \\
alpaca-0255 & -- & UF-DPO & SFT advises deleting Anki's data files (which would erase the user's
cards) and cites a non-existent support team. \\
alpaca-0269 & -- & UF-DPO & SFT fails the problem and falls into a repetition loop. \\
\multicolumn{4}{@{}l}{\emph{UltraFeedback: DPO regresses}} \\
alpaca-0142 & -- & SFT & SFT correctly says there is not enough information; the longer DPO answer
invents details (``15\% cashback'') and does nonsense arithmetic. \\
alpaca-0089 & -- & SFT & Both describe racquetball wrongly; DPO's longer answer adds a bizarre
claim that players hit each other with racquets. \\
\bottomrule
\end{tabularx}
\end{table}

The UF regressions show the mechanism suggested by the length statistics: the extra length is
often filled with invented specifics. The UF wins mostly come from avoiding SFT failure modes
(harmful advice, repetition loops) rather than from more informative answers. In HH, DPO wins when
caution or careful constraint-following is the right behaviour and loses when the task simply
rewards a fuller answer.

% ===========================================================================
\section{Limitations and Reproducibility}
% ===========================================================================
\textbf{Limitations.} (i) A single LLM judge with its own biases; order bias is mitigated by
randomisation but not measured. (ii) 300 prompts per set give 95\% intervals of roughly
$\pm$5--6 points on a win rate near 50\%. (iii) One seed per run, so run-to-run variance is unknown.
(iv) Each DPO run is a single epoch and HH metrics were still improving. (v) The 512-token limit
truncates many UltraFeedback responses, and the SFT length filter biases towards short answers.
(vi) Only single-turn data are used (29\% of the HH pool). (vii) A 0.5B model hallucinates
regardless of alignment, which limits absolute quality. (viii) The persona evaluation has only 40
prompts, and 5 of them could not be judged. (ix) 8 (HH) and 13 (UF) verdicts were unparseable and are
excluded; this is small relative to $n=300$ but may not be random with respect to response content.

\textbf{Reproducibility.} Kaggle notebooks: \texttt{pa2-01-prepare-data}, \texttt{pa2-02-sft},
\texttt{pa2-03-dpo-hh}, \texttt{pa2-04-dpo-uf}, \texttt{pa2-05-generate}, \texttt{pa2-07-persona},
\texttt{pa2-06-judge}; each run's \texttt{config\_resolved.yaml} records the exact values used.
\begin{center}\footnotesize
\begin{tabular}{@{}ll@{}}
\toprule
Item & Path in the repository \\
\midrule
Configurations & \texttt{configs/\{sft,dpo\_hh,dpo\_uf,persona\_lora,generate\}.yaml} \\
Run metrics & \texttt{report/runs/\{sft,dpo\_hh,dpo\_uf,persona\_lora\}/} (metrics, best step, resolved config) \\
Data statistics & \texttt{report/data\_stats.json} \\
Evaluation prompts & \texttt{eval\_sets/\{hh\_rlhf\_eval\_300\_seed42,alpaca\_eval\_300\_seed42,persona\_style\_prompts\}.json} \\
Generations & \texttt{report/generations/*.json} (+ \texttt{*.meta.json} with decoding settings) \\
Blind A/B pairs & \texttt{report/pairs/*.jsonl} \\
Judgments and summaries & \texttt{report/judgments/*.jsonl}, \texttt{*\_failures.jsonl}, \texttt{*.summary.\{json,md\}} \\
Figures and tables & \texttt{report/figures/}, \texttt{report/tables/} (\texttt{scripts/make\_report\_assets.py}) \\
\bottomrule
\end{tabular}
\end{center}
Commands: \texttt{scripts/prepare\_data.py}, \texttt{train\_sft.py}, \texttt{train\_dpo.py --config
configs/dpo\_\{hh,uf\}.yaml}, \texttt{train\_persona\_lora.py}, \texttt{generate.py},
\texttt{build\_pairs.py}, \texttt{judge\_pairwise\_kaggle.py}, \texttt{summarize\_judgments.py}.

\end{document}
